\documentclass[11pt,reqno]{amsart}

%==============================================================
% An Energy-First Reorganization of Stochastic Calculus
% (Framework paper; IBM treated in the companion paper.)
%==============================================================

\usepackage[utf8]{inputenc}
\usepackage[T1]{fontenc}
\usepackage{lmodern}
\usepackage[margin=1.15in]{geometry}
\usepackage{amsmath,amssymb,amsthm,mathtools}
\usepackage{mathrsfs}
\usepackage{microtype}
\usepackage{hyperref}
\usepackage{cleveref}
% --- subsection style: centered, small caps, unnumbered visually, no
% trailing period. Internal counter is preserved so \cref still works.
% (Direct redefinition; titlesec conflicts with amsart's section internals.)
\makeatletter
\renewcommand\subsection{\@startsection{subsection}{2}%
  \z@{.7\linespacing}{.5\linespacing}%
  {\normalfont\centering\scshape}}
\makeatother

% --- theorem environments
\theoremstyle{plain}
\newtheorem{theorem}{Theorem}[section]
\newtheorem{proposition}[theorem]{Proposition}
\newtheorem{lemma}[theorem]{Lemma}
\newtheorem{corollary}[theorem]{Corollary}

\theoremstyle{definition}
\newtheorem{definition}[theorem]{Definition}
\newtheorem{example}[theorem]{Example}
\newtheorem{assumption}[theorem]{Assumption}
\newtheorem{conjecture}[theorem]{Conjecture}
\newtheorem{program}[theorem]{Program}

\theoremstyle{remark}
\newtheorem{remark}[theorem]{Remark}
\newtheorem*{sketch}{Sketch}

% --- macros
\newcommand{\R}{\mathbb{R}}
\newcommand{\N}{\mathbb{N}}
\newcommand{\E}{\mathbb{E}}
\newcommand{\Prob}{\mathbb{P}}
\newcommand{\Hilb}{\mathcal{H}}
\newcommand{\Dom}{\mathrm{Dom}}
\newcommand{\Ran}{\mathrm{Ran}}
\newcommand{\Gra}{\mathrm{Gra}}
\newcommand{\Energy}{\mathcal{E}}
\newcommand{\Source}{\mathsf{S}}
\newcommand{\Obs}{\mathsf{O}}
\renewcommand{\iota}{\mathfrak{i}}
\newcommand{\IBM}{Z}
\newcommand{\Carre}{\Gamma}
\DeclareMathOperator{\Var}{Var}
\DeclareMathOperator{\Cov}{Cov}

% Make the author block render in natural case (amsart's default
% \@setauthors wraps the author list in \MakeUppercase).
\makeatletter
\renewcommand{\@setauthors}{%
  \begingroup
  \trivlist
  \centering\footnotesize\@topsep30\p@\relax
  \advance\@topsep by -\baselineskip
  \item\relax
  \andify\authors
  \def\\{\protect\linebreak}%
  \authors\relax
  \endtrivlist
  \endgroup
}
\makeatother

\title[Energy-First Stochastic Calculus]{An Energy-First Reorganization of Stochastic Calculus:\\
       Operator-Derivative Calculus on Cylindrical Noises}
\author{Ramiro Fontes\\ Quijotic Research\\ \texttt{ramirofontes@gmail.com}\\ \texttt{www.quijotic.com}}
\date{\today}

\begin{document}

\begin{abstract}
We reverse the conventional order of constructing stochastic calculus.
The classical sequence builds the stochastic integral first, defines
the derivative by adjointness, and recovers the It\^o formula at the
end. We take the operator derivative
$D\colon\Dom(D)\subseteq\Hilb\to\Energy$ on a Hilbert energy space as
the primitive datum and define the divergence by adjunction,
\begin{equation}\label{eq:abstract-reversal}
  \delta\;\coloneqq\;D^{*}.
\end{equation}

The consequence of this reversal is that It\^o's formula becomes a
parametric family indexed by the \emph{observable correction order}
$q$ (equivalently the spatial order of the PDE shadow), with the
underlying operator datum carrying its own integer \emph{operator
rank} $r$. In this framework It\^o's formula reads, parametrically
in $q$,
\begin{equation}\label{eq:abstract-ito-hierarchy}
  f(U_t)-f(U_0)\;=\;\delta\bigl(B_f(t)\bigr)\;+\;\mathrm{Corr}_f^{[q]}(t),
  \qquad\mathrm{Corr}_f^{[q]}(t)\in\bigl(\Energy^{[q]}\bigr)^{*}.
\end{equation}
For $q=2$ the dual $(\Energy^{[2]})^{*}$ is realised by a.s.-finite
real random variables paired against $f''(X)\,d\Carre^{D}_{X}$,
recovering the classical $L^{2}(\Omega)$ correction; for $q\ge 4$ the
correction genuinely lives in a distributional dual $\Energy^{[q],-}$
via the rigged hierarchy of \cref{prop:quartic-dual} (see also
\cref{rem:schematic-correction}).
This identity is the framework's organizing form: when the source
noise admits an It\^o expansion and the source integrand lifts through
the observable inclusion, the conditional transfer principle produces
an observable-level identity of this shape, with $q$ determined by the
order of the observable correction energy space. The parametric
identity is developed as the framework's central organizing thesis in
\cref{sec:energy-order} and instantiated for specific noise structures
in the companion papers. The classical quadratic-variation correction
$\tfrac{1}{2}\!\int_{0}^{t}f''(X_s)\,d\langle X\rangle_{s}$ is the
$q=2$ case, corresponding to operator rank $r=1$. The IBM construction
in the companion paper~\cite{IBMPaper} is, conditional on the companion
paper, the $q=4$ case (cf.~\cref{cor:k4-witness}), corresponding to
operator rank $r=2$, for iterated Brownian motion
$\IBM_t=X_{|Y_t|}$. In the even-order tensor cases considered in this
paper, $q=2r$.

The abstract framework that produces
\eqref{eq:abstract-ito-hierarchy} consists of three theorems: the
compression construction (\cref{thm:compression}), which closes the
pullback $D_\Obs F=D_\Source(\iota F)$ of a source derivative through
an observable sub-$\sigma$-algebra; the projected divergence theorem
(\cref{thm:proj-div}), which reads off the intrinsic divergence
$\delta_\Obs u=\E[\delta_\Source(Ju)\mid\mathcal F_\Obs]$ as a
conditional expectation of the source divergence; and the conditional
operator-It\^o transfer principle (\cref{thm:transfer}), which lifts
a source-level It\^o expansion to the observable level.
\end{abstract}

\maketitle

\section{Introduction}\label{sec:intro}

The classical construction of stochastic calculus, in nearly every
textbook treatment, proceeds in a fixed order: fix a Brownian motion
$B$ on a probability space; construct the stochastic (It\^o or
Skorohod) integral $\delta$ on a class of integrands by an isometry
argument; define the Malliavin derivative $D$ either by chaos
expansion or as the adjoint $D=\delta^{*}$; finally derive the
It\^o/Tanaka formulae and identify the PDE shadow of the process as
the action of the infinitesimal generator on smooth functions. For
Brownian motion this ordering is harmless --- the Gaussian structure
forces every reasonable derivative/divergence pair to coincide, the
energy space is $L^{2}(\R_{+})$ with the Lebesgue inner product, and
the PDE shadow is the heat equation. The ordering nevertheless hides
a structural choice. \emph{Three} pieces of data are in play: an
energy space $\Energy$, an operator derivative
$D\colon L^{2}(\Omega)\to\Energy$, and a divergence
$\delta\colon\Energy\to L^{2}(\Omega)$. In the Brownian regime any two
determine the third. Outside that regime --- fractional Brownian
motion, Volterra and ambit processes, stable L\'evy processes,
stochastic volatility models, and composite (subordinated,
time-changed, iterated) processes such as iterated Brownian motion ---
the choice of $\Energy$ becomes the primary modelling decision, and
postulating $\delta$ first is a technical convenience rather than the
conceptually correct order.

The companion paper~\cite{Fontes-OperatorDerivative} (\emph{The
Operator Derivative in Continuous Stochastic Calculus}) develops a
divergence-first operator framework: it takes the stochastic integral
$\delta_X$ as primitive on a Hilbert energy space $\mathcal H_X$ and
defines the operator derivative by adjointness
$D_X\coloneqq\delta_X^{*}$ on its natural adjoint domain. The
representation layer (Clark--Ocone, GKW projection, variance
identity) is shown there to follow from adjointness alone; the
calculus layer (Leibniz, chain rule, It\^o formula) requires a
cylindrical reduction with process-dependent representers. The
present paper inverts the order of that construction. The primitive
object is the pair $(\Energy, D)$; the divergence is
\emph{defined} as $\delta\coloneqq D^{*}$, and the entire calculus is
rebuilt downstream of this single reversal. The term ``operator
derivative'' is retained throughout for $D$, with the understanding
that the framing is now energy-first --- $\Energy$ and $D$ are the
data, $\delta$ is derived --- rather than divergence-first.

For a composite process in which an observable $\sigma$-algebra
$\mathcal F_\Obs$ is strictly contained in a richer source
$\sigma$-algebra $\mathcal F_\Source$, the reversal goes one step
further. A source energy datum
$(\Hilb_\Source,\Energy_\Source,D_\Source)$ induces an intrinsic
observable energy datum $(\Hilb_\Obs,\Energy_\Obs,D_\Obs)$ by
\emph{compression} --- restricting $D_\Source$ to functionals of the
observable $\sigma$-algebra and closing on the resulting graph. The
intrinsic divergence is the observable conditional expectation of the
source divergence, given by the projected
formula~\eqref{eq:proj-div} on the lifted source-divergence
domain; the compression and projected divergence are
\cref{thm:compression,thm:proj-div} of \cref{sec:compat}. Layered on
top of compression is a conditional operator-It\^o transfer
principle (\cref{thm:transfer}): a source-level It\^o identity
$f(U_{t})-f(U_{0})=\delta_\Source(A_{f}(t))+R_{f}(t)$ converts into
an observable-level identity
$f(U_{t})-f(U_{0})=\delta_\Obs(B_{f}(t))+P_\Obs R_{f}(t)$ whenever
the source integrand $A_{f}(t)$ lifts through the observable
inclusion $J$. The practical content is a separation of duties: the
source side supplies the It\^o expansion; the compression
construction supplies the observable lift; the correction term is
the observable projection of the source remainder. The transfer is
what makes the energy-first framework operational in concrete
examples without forcing the observable to be a semimartingale in
its own filtration.

A second structural payoff of the reversal is a classification of
It\^o corrections by their \emph{observable correction order} $q$
(equivalently the order of the PDE shadow), with a separate integer
\emph{operator rank} $r$ recording the order of the operator-derivative
datum that realises the correction. The central object of the
classification is the parametric It\^o identity
\eqref{eq:abstract-ito-hierarchy}: under the hypotheses of the
transfer theorem (\cref{thm:transfer}) --- a source-level It\^o
expansion of the observable and a source integrand that lifts through
the observable inclusion --- the transfer principle produces an
It\^o-type expansion of the observable whose centred part is a
divergence and whose correction lives in the dual of an order-$q$
observable correction energy space, with the order $q$ determined by
the noise composition. The framework does not claim such an expansion
unconditionally for every observable; the source-expansion and lift
hypotheses are genuine compatibility requirements, verified case by
case in the companion papers. Classical semimartingale It\^o
calculus organises corrections by \emph{quadratic variation}: the
correction in
$f(X_{t})-f(X_{0})=M_{f}(t)+\tfrac{1}{2}\!\int_{0}^{t}f''(X_s)\,
d\langle X\rangle_{s}$ is encoded by the second-order bracket
$d\langle X\rangle$, which is the $q=2$ instance of
\eqref{eq:abstract-ito-hierarchy} (and is generated by a single
operator derivative, operator rank $r=1$). The energy-first viewpoint
replaces the quadratic-variation organising principle with the
\emph{correction order} of the observable: the case $q=2$, $r=1$,
$\Carre(F,F)=\langle DF,DF\rangle_\Energy$ recovers
$d\langle X\rangle$, but higher-order observable correction energies
$\Energy^{[q]}$ --- generated by higher-rank operator-derivative data
$D^{(r)}$ via the relation $q=2r$ in the even-order tensor cases ---
encode It\^o corrections whose PDE shadows are differential operators
of order $q>2$. The non-semimartingale obstruction --- absence of a
usable $d\langle X\rangle$ --- is then not an obstacle but a coordinate
artefact: the observable simply lives at higher correction order than
the classical $q=2$ bracket assumes. This thesis is formulated and
developed in \cref{sec:energy-order}.

The framework was developed against the concrete obstruction posed by
iterated Brownian motion $\IBM_{t}=X_{|Y_{t}|}$, with two-sided
Brownian scenery $X$ and an independent Brownian clock $Y$. In its
natural filtration $\sigma(\IBM_s:s\le t)$, $\IBM$ is not a
semimartingale and does not carry the finite predictable bracket
structure required by the classical semimartingale It\^o formula, so
the classical construction of an It\^o formula fails outright. The companion
paper~\cite{IBMPaper} carries out the energy-first construction for
this example: it applies the compression and the operator-It\^o
transfer to the source $(X,Y)$, identifies the fourth-order
observable energy that organises the resulting It\^o correction,
proves a mollified pathwise It\^o formula, and isolates the
$L^{2}(\Omega)$ obstruction to the unmollified pathwise correction.
IBM is the canonical fourth-order witness of the energy-order
classification proposed here. The present paper develops only the
abstract framework; no IBM-specific calculations beyond the
identification of IBM as motivating example appear here. The
source-resolved chaos atlas paper~\cite{Fontes-IBMchaos} supplies the
$(m,n)$-resolved chaos kernels used in~\cite{IBMPaper};
\cite{Fontes-IBMnumerical} validates the bidegree variance budgets
numerically; \cite{Fontes-PDEgenerators} treats PDE shadows of energy
structures generally. None of those papers are used in the present
framework presentation.

\subsection*{Main results}

\Cref{thm:compression} (compression construction): a source datum
$(\Hilb_\Source,\Energy_\Source,D_\Source)$ with observable
$\sigma$-algebra $\mathcal F_\Obs\subseteq\mathcal F_\Source$ induces,
on the maximal pullback domain
$\mathcal C_\Obs=\{F\in\Hilb_\Obs:\iota F\in\Dom(D_\Source)\}$, a
closed first-order intrinsic observable derivative
$D_\Obs F=D_\Source(\iota F)$ valued in the observable energy
$\Energy_\Obs=\overline{D_\Source(\iota \mathcal C_\Obs)}^{\Energy_\Source}$.
\Cref{thm:proj-div} (projected divergence): the adjoint
$\delta_\Obs=D_\Obs^{*}$ admits, on the lifted source-divergence
domain, the conditional-expectation formula
$\delta_\Obs u=\E[\delta_\Source(Ju)\mid\mathcal F_\Obs]$.
\Cref{thm:transfer} (conditional operator-It\^o transfer): a
source-level It\^o expansion lifts to an observable-level identity
whenever the source integrand lies in $J(\Energy_\Obs)$, with the
observable correction obtained as the observable projection of the
source remainder.

\subsection*{Contents}

\Cref{sec:operator} sets up the three-level hierarchy of It\^o
formulae --- observable It\^o identity, pathwise lift, energy-first /
operator lift --- and the energy datum that organises Level~3. The
compression construction and the projected divergence formula are
proved in \cref{sec:compat}; the operator-It\^o transfer principle in
\cref{sec:transfer}. \Cref{sec:energy-order} formulates the
energy-order classification as the central organising thesis of the
reorganization, with the second-order quadratic-variation case
recovered as a corollary. \Cref{sec:outlook} closes with a forward
pointer to the IBM construction of~\cite{IBMPaper}.

\section{The three levels of an It\^o formula}\label{sec:operator}

We work throughout on a probability space $(\Omega,\mathcal F,\Prob)$
and write $L^2(\Omega)=L^2(\Omega,\mathcal F,\Prob)$. This section sets
up the three nested levels at which the framework reads an It\^o
formula and introduces the energy datum that organises the
energy-first / operator layer.

The framework is organised around the principle that an It\^o formula
should be read at three nested levels of structure: a baseline
\emph{observable} layer, a \emph{pathwise} lift, and an
\emph{energy-first / operator} lift. The classical semimartingale
It\^o formula is the quadratic-energy case of the third level; it is
one specific lift, not the definition.

\begin{definition}[Observable It\^o identity]\label{def:obs-ito-identity}
Let $\mathcal T\subseteq C(\R)$ be a class of test functions and
$X=(X_t)_{t\ge 0}$ an $\R$-valued process with
$\E|f(X_t)|<\infty$ for every $f\in\mathcal T$ and $t\ge 0$. We
say that $X$ admits an \emph{observable It\^o identity} on
$\mathcal T$ if, for each $f\in\mathcal T$, the map
$t\mapsto\E[f(X_t)]$ is locally absolutely continuous on
$(0,\infty)$ and there exist a family of linear operators
$\mathcal G_t$ on $\mathcal T$ and a family of linear functionals
$\mathcal B_t$ on $\mathcal T$ such that
\begin{equation}\label{eq:obs-ito-identity}
  \frac{d}{dt}\,\E[f(X_t)]
   \;=\;\E\bigl[(\mathcal G_t f)(X_t)\bigr]
   \;+\;\mathcal B_t f
\end{equation}
in the sense of distributions in $t$, or pointwise wherever the
left-hand-side derivative exists.
\end{definition}

We call $\mathcal G_t$ the \emph{observable correction operator}
and $\mathcal B_t$ the \emph{boundary/source term}. For Brownian
motion, $\mathcal G=\tfrac12\partial_x^{2}$ and $\mathcal B=0$; for
a general Markov process the observable correction operator is the
infinitesimal generator. We do not require $X$ to be Markovian, nor
to be a semimartingale: \cref{def:obs-ito-identity} is purely an
evolution identity for test-function expectations, and assumes no
filtration structure on $X$.

On a given class $\mathcal T$, the hierarchy reads:
\begin{itemize}
\item[\textbf{Level 1.}] \emph{Observable It\^o identity}
  \eqref{eq:obs-ito-identity}: an evolution identity for
  $t\mapsto\E[f(X_t)]$. This is the weakest level; it depends only
  on the one-dimensional time marginals of $X$ paired with
  $\mathcal T$.
\item[\textbf{Level 2.}] \emph{Pathwise lift}: a decomposition
  \begin{equation}\label{eq:pathwise-lift}
    f(X_t)-f(X_0)\;=\;M_f(t)\;+\;C_f(t),
    \qquad \E[M_f(t)]=0,
  \end{equation}
  where $C_f$ is a process whose time-derivative expectation
  reproduces the right-hand side of \eqref{eq:obs-ito-identity}
  (so that integrating in time and taking expectation recovers
  Level~1), and $M_f$ is a centred fluctuation.
\item[\textbf{Level 3.}] \emph{Energy-first / operator lift}: a
  further lift of \eqref{eq:pathwise-lift} in which $M_f$ is
  represented by a divergence/integral $\delta(U_f)$ for an energy
  datum $(\Hilb,\Energy,D)$ in the sense of \cref{def:energy-datum}
  below, and $C_f$ is represented by an observable-energy
  contraction associated with $D$ or a higher-order analogue
  of $D$.
\end{itemize}

Each lift in this ladder is additional structure, not an automatic
consequence of the level below it. A Level-1 observable It\^o
identity does not by itself produce a Level-2 pathwise lift: such a
lift requires $C_f$ to exist
pathwise in a suitable topology, which is a structural property of
$X$ relative to $\mathcal T$ and may fail when $X$ is not a
semimartingale or when $\mathcal G_t$ is of higher order than two.
Similarly, a Level-2 lift is not automatically Level-3: realising
$M_f$ as $\delta(U_f)$ and $C_f$ as an energy contraction requires
an energy datum whose divergence/contraction reproduces the
pathwise decomposition. The classical semimartingale It\^o formula
sits at Level~3 with a bracket-compatible operator-energy datum
$(\Hilb,\Energy,D)$ and the quadratic-energy contraction
$C_f(t)=\tfrac12\int_0^t f''(X_s)\,d\langle X\rangle_s$: it is one
specific energy-first lift, not the definition of an It\^o formula.
In the canonical Brownian/Wiener specialisation, the operator
derivative $D$ may be taken to be the standard Malliavin derivative on
$\Energy=L^{2}(\R_{+})$, but for a general semimartingale the
recovery only requires a bracket-compatible operator-energy datum in
the sense of \cref{thm:k2-recovery}.

The framework does not constrain the order of the observable
correction operator $\mathcal G_t$. Brownian motion is a second-order
observable It\^o process. A fourth-order observable It\^o process is
exhibited in the worked example of~\cite{IBMPaper}: iterated Brownian
motion $\IBM_t=X_{|Y_t|}$ has $\mathcal G_t=\tfrac18\partial_x^{4}$
plus a boundary term, and is therefore a fourth-order observable
It\^o process even though it is not a semimartingale in its own
filtration. The energy datum that organises the corresponding
Level-3 lift is not a quadratic-energy / Malliavin datum but a
genuinely fourth-order observable energy, and this is the
energy-order classification we formulate in
\cref{sec:energy-order} as the central organising thesis of the
framework.

\subsection*{Energy datum and the operator/divergence layer}

We now state the structural data underlying Level~3. The rest of
this section sets up the energy datum, the carr\'e du champ, and
the three-layer (R)/(P)/(I) decomposition internal to Level~3.

\begin{definition}[Energy datum]\label{def:energy-datum}
An \emph{energy datum} is a triple $(\Hilb,\Energy, D)$ where
$\Hilb\subseteq L^2(\Omega)$ is a closed subspace,
$\Energy$ is a separable Hilbert space, and
\[
  D\colon \Dom(D)\subseteq\Hilb \longrightarrow \Energy
\]
is a closed, densely defined linear operator. The
\emph{divergence} is the Hilbert-space adjoint
$\delta\coloneqq D^{*}\colon\Dom(\delta)\subseteq\Energy\to\Hilb$.
\end{definition}

The closedness assumption is what makes adjunction well-posed: by the
von Neumann adjoint theorem, $\delta=D^{*}$ is then also closed and
densely defined.

\begin{definition}[Carr\'e du champ]\label{def:carre}
On $\Dom(D)\times\Dom(D)$ define the bilinear form
$\Carre(F,G)\coloneqq\langle DF,\,DG\rangle_\Energy$.
We call $\Carre$ the \emph{energy bracket} associated to the datum.
\end{definition}

In the canonical Brownian/Wiener specialisation
$\Hilb=L^2(\Omega)$, $\Energy=L^2(\R_+)$, with $D$ taken to be the
standard Malliavin derivative, the bracket $\Carre$ reduces to the
standard Malliavin matrix entry; for general operator-energy data
$(\Hilb,\Energy,D)$ the present formalism does not require this
specialisation.

Once the datum is fixed, three layers are constructed in sequence.

\medskip
\noindent\textbf{(R) Representation layer.}
Given $(\Hilb,\Energy,D)$, the integration-by-parts formula
\begin{equation}\label{eq:ibp}
  \E\bigl[F\,\delta u\bigr]
  \;=\;\langle DF,\,u\rangle_{\Energy},\qquad
  F\in\Dom(D),\;u\in\Dom(\delta),
\end{equation}
holds by definition of adjunction. \emph{Existence of a chaos
expansion} adapted to $D$ is then an independent assumption on
$(\Hilb,\Energy,D)$, e.g.~that $D$ admits a Wiener chaos
diagonalisation.

\medskip
\noindent\textbf{(P) Product/chain layer.}
A \emph{product rule} for $D$, of the form
$D(FG)=F\,DG+G\,DF$ on a suitable algebra, and a \emph{chain rule}
$D(f(F))=f'(F)\,DF$ on smooth $f$, are not consequences of
\cref{def:energy-datum}. They have to be imposed as further hypotheses
on the datum; when they hold, we call the datum a \emph{calculus datum}.

\medskip
\noindent\textbf{(I) It\^o layer.}
Once (R) and (P) are in place, an It\^o-type formula
\begin{equation}\label{eq:operator-ito}
  f(X_t)-f(X_0) \;=\; \delta\bigl(U_f(X_\bullet)\bigr) + C_f(X_\bullet),
\end{equation}
with $U_f$ a representation integrand built from $f',DX$ and
$C_f$ an \emph{energy correction} built from $f''$ and $\Carre(X,X)$,
becomes a candidate identity. We do not assert that
\cref{eq:operator-ito} holds for every calculus datum; we record it as
the scheme that must be verified in each instance.

\begin{remark}[Conditional, not automatic]\label{rem:conditional}
The energy-first framework does not produce \cref{eq:operator-ito}
for free. What it does produce is a sharp separation between the
two ingredients of an It\^o formula: a divergence representative
$U_f$ in $\Energy$, and a correction term whose form is dictated by
$\Carre$. For Gaussian/Brownian data the correction reduces to the
standard $\tfrac12 f''(X)\,d\langle X\rangle$; for non-Gaussian or
composite data it need not.
\end{remark}

\begin{remark}[Two-layer separation, inherited]\label{rem:two-layer}
The representation/calculus separation
of~\cite{Fontes-OperatorDerivative} survives the reversal: the
integration-by-parts formula~\eqref{eq:ibp} and any chaos
representation adapted to $D$ form the representation layer, while a
chain/product rule for $D$ is the additional input needed for the
calculus layer. Such a chain rule is equivalent, in the sense
of~\cite[Thm.~B]{Fontes-OperatorDerivative}, to a \emph{cylindrical
reduction} --- the existence of a dense algebra of test functionals
$f(\xi_1,\dots,\xi_n)$ and representers $\kappa_i\in L^2(\Omega;\Energy)$
through which the integration-by-parts formula factors. The
classical Malliavin calculus is the special case in which the
representers are the deterministic Cameron--Martin elements; in the
present setting it is recovered as one cylindrical reduction among
others, not as the definition of $D$. The difference between the two
orderings is one of \emph{starting datum}: in the energy-first
ordering, both layers are read off $(\Energy, D)$ rather than off
$(\mathcal H_X, \delta_X)$.
\end{remark}

\section{The compression construction and the projected divergence formula}\label{sec:compat}

This section contains the central abstract content of the framework.
Given a closed densely defined source derivative $D_\Source$ and an
observable $\sigma$-algebra $\mathcal F_\Obs\subseteq\mathcal F_\Source$,
we construct a closed first-order intrinsic derivative $D_\Obs$ and its
adjoint $\delta_\Obs$. On the subspace of $u\in\Energy_\Obs$ for which
$Ju\in\Dom(\delta_\Source)$, the intrinsic divergence is given by the
projected formula $\delta_\Obs u=P_\Obs\delta_\Source(Ju)$; we do not
claim this characterises all of $\Dom(\delta_\Obs)$. Exact
intertwining $\iota\,\delta_\Obs=\delta_\Source J$ is recovered as a
special case.

\subsection*{Setup}

Let $\mathcal F_\Source\supseteq\mathcal F_\Obs$ be sub-$\sigma$-algebras
of $\mathcal F$. Set
\[
  \Hilb_\Source\coloneqq L^2(\Omega,\mathcal F_\Source,\Prob),\qquad
  \Hilb_\Obs\coloneqq L^2(\Omega,\mathcal F_\Obs,\Prob),
\]
and let
\[
  \iota\colon \Hilb_\Obs\hookrightarrow\Hilb_\Source
\]
denote the canonical isometric inclusion. Its Hilbert-space adjoint is
the conditional expectation onto $\mathcal F_\Obs$,
\[
  P_\Obs\coloneqq\iota^{*}=\E[\,\cdot\mid\mathcal F_\Obs]\colon
  \Hilb_\Source\to\Hilb_\Obs.
\]
Let $(\Hilb_\Source,\Energy_\Source,D_\Source)$ be a source energy datum
in the sense of \cref{def:energy-datum}, with divergence
$\delta_\Source=D_\Source^{*}$.

\subsection*{The compression datum}

\begin{definition}[Observable smooth core]\label{def:obs-core}
The \emph{observable smooth core} induced by $D_\Source$ is
\[
  \mathcal C_\Obs
  \;\coloneqq\;
  \bigl\{\,F\in\Hilb_\Obs : \iota F\in\Dom(D_\Source)\,\bigr\}.
\]
\end{definition}

\begin{assumption}[Density of the observable core]\label{ass:dense}
$\mathcal C_\Obs$ is dense in $\Hilb_\Obs$.
\end{assumption}

\begin{remark}[Density routes]\label{rem:density-routes}
In the canonical cases (Wiener, Poisson, product Gaussian source
noises), density of $\mathcal C_\Obs$ in $\Hilb_\Obs$ follows from
standard polynomial-functional or cylindrical-functional density
theorems. For non-classical noise structures (Tsirelson--Vershik
black noises, sticky-flow noises, composite source pairs), density
is a genuinely structural condition that is verified case by case in
the companion papers (see~\cite{IBMPaper} for the iterated Brownian
motion source pair $(X,Y)$; \cite{StickyFlowNoiseLift} for the LJR
sticky-flow case). The framework paper assumes density as a
structural input; the verifications belong to the analytical
companion papers.
\end{remark}

\Cref{ass:dense} is a structural input. Its verification depends on
$\mathcal F_\Obs\hookrightarrow\mathcal F_\Source$ and on $D_\Source$.
For abstract Wiener spaces with a smooth source datum it is typically
inherited from smooth cylindrical functionals; in concrete examples
where $\mathcal F_\Obs$ is the $\sigma$-algebra of a composite process
the verification is generally non-trivial and is among the analytic
tasks left for the application paper (see, e.g.,~\cite{IBMPaper} for
the IBM specialisation).

\begin{definition}[Compression operator and observable energy]
\label{def:compression}
For $F\in\mathcal C_\Obs$, set
\[
  D_\Obs^{0} F \;\coloneqq\; D_\Source(\iota F) \;\in\; \Energy_\Source.
\]
Define the observable energy space
\[
  \Energy_\Obs
  \;\coloneqq\;
  \overline{D_\Obs^{0}\,\mathcal C_\Obs}^{\,\Energy_\Source}
  \;\subseteq\;\Energy_\Source,
\]
a closed subspace of $\Energy_\Source$, and let
$J\colon\Energy_\Obs\hookrightarrow\Energy_\Source$ denote the
isometric inclusion. Regard $D_\Obs^{0}$ henceforth as a linear
operator $\mathcal C_\Obs\to\Energy_\Obs$.
\end{definition}

\subsection*{Closability and the intrinsic derivative}

\begin{theorem}[Observable compression derivative]\label{thm:compression}
Under \cref{ass:dense}, the operator
$D_\Obs^{0}\colon\mathcal C_\Obs\subseteq\Hilb_\Obs\to\Energy_\Obs$ is
closed on the maximal pullback domain $\mathcal C_\Obs$. We denote it
$D_\Obs$ and call it the \emph{minimal observable derivative} induced by
$D_\Source$. Consequently, any restriction
$D_\Obs^{0}|_{\mathcal C'}$ to a subcore
$\mathcal C'\subseteq\mathcal C_\Obs$ that is dense in $\Hilb_\Obs$ is
closable, and the closure of any such restriction is contained in
$D_\Obs$.
\end{theorem}

\begin{proof}
Let $F_n\in\mathcal C_\Obs$ with $F_n\to F$ in $\Hilb_\Obs$ and
$D_\Obs^{0}F_n\to h$ in $\Energy_\Obs$. Since $\iota$ is an isometric
inclusion, $\iota F_n\to\iota F$ in $\Hilb_\Source$; since $J$ is an
isometric inclusion,
$J\,D_\Obs^{0}F_n=D_\Source(\iota F_n)\to Jh$ in $\Energy_\Source$.
Closedness of $D_\Source$ yields $\iota F\in\Dom(D_\Source)$, hence
$F\in\mathcal C_\Obs$, and $D_\Source(\iota F)=Jh$, whence
$D_\Obs^{0}F=h$. So the graph of $D_\Obs^{0}$ on $\mathcal C_\Obs$ is
closed.

The second statement is a direct consequence: if
$\mathcal C'\subseteq\mathcal C_\Obs$ is dense in $\Hilb_\Obs$, the same
graph-closure argument applied to sequences in $\mathcal C'$ converging
in $\Hilb_\Obs\times\Energy_\Obs$ to a graph point lands in
$\mathcal C_\Obs$ and in the graph of $D_\Obs$. In particular, the
closure of $D_\Obs^{0}|_{\mathcal C'}$ exists as a closed operator and
is contained in $D_\Obs$.
\end{proof}

\begin{remark}[Role of the density hypothesis]\label{rem:density-role}
The density hypothesis \cref{ass:dense} is used in the framework
downstream: it is what makes $D_\Obs$ densely defined, and hence
admissible as the derivative of a compression energy datum in the
sense of \cref{def:energy-datum}. The closedness conclusion of
\cref{thm:compression} uses only closedness of $D_\Source$ and the
isometricity of $\iota$ and $J$.
\end{remark}

The intrinsic divergence is now defined by adjunction.

\begin{definition}[Intrinsic divergence]\label{def:int-div}
\[
  \delta_\Obs\coloneqq D_\Obs^{*}\colon\Dom(\delta_\Obs)\subseteq\Energy_\Obs
  \longrightarrow\Hilb_\Obs.
\]
\end{definition}

\subsection*{The projected divergence formula}

\begin{theorem}[Projected divergence formula]\label{thm:proj-div}
Let $u\in\Energy_\Obs$ satisfy $Ju\in\Dom(\delta_\Source)$. Then
$u\in\Dom(\delta_\Obs)$ and
\begin{equation}\label{eq:proj-div}
  \delta_\Obs u \;=\; P_\Obs\,\delta_\Source(Ju).
\end{equation}
\end{theorem}

\begin{proof}
Let $F\in\mathcal C_\Obs$. By \cref{thm:compression} and the inclusion
isometry $J$,
\[
  \langle D_\Obs F,\,u\rangle_{\Energy_\Obs}
  \;=\;\langle J\,D_\Obs F,\,Ju\rangle_{\Energy_\Source}
  \;=\;\langle D_\Source(\iota F),\,Ju\rangle_{\Energy_\Source}.
\]
By source integration-by-parts (\cref{eq:ibp} applied to $D_\Source$)
and the hypothesis $Ju\in\Dom(\delta_\Source)$,
\[
  \langle D_\Source(\iota F),\,Ju\rangle_{\Energy_\Source}
  \;=\;\langle\iota F,\,\delta_\Source(Ju)\rangle_{\Hilb_\Source}.
\]
Applying $\iota^{*}=P_\Obs$,
\[
  \langle\iota F,\,\delta_\Source(Ju)\rangle_{\Hilb_\Source}
  \;=\;\langle F,\,P_\Obs\,\delta_\Source(Ju)\rangle_{\Hilb_\Obs}.
\]
Chaining,
\begin{equation}\label{eq:chain}
  \langle D_\Obs F,\,u\rangle_{\Energy_\Obs}
  \;=\;\langle F,\,P_\Obs\,\delta_\Source(Ju)\rangle_{\Hilb_\Obs}
  \qquad\text{for all }F\in\mathcal C_\Obs.
\end{equation}
The right-hand side of \cref{eq:chain} is a bounded linear functional of
$F\in\Hilb_\Obs$, with bound
$\|P_\Obs\,\delta_\Source(Ju)\|_{\Hilb_\Obs}$. Since
$\Dom(D_\Obs)=\mathcal C_\Obs$ by \cref{thm:compression}, the identity
holds on all of $\Dom(D_\Obs)$. Hence $u\in\Dom(D_\Obs^{*})=
\Dom(\delta_\Obs)$ and $\delta_\Obs u=P_\Obs\,\delta_\Source(Ju)$.
\end{proof}

\begin{remark}[Automaticity of the $L^{2}$ hypothesis]
\label{rem:l2-auto}
In more abstract versions of this construction one finds the additional
hypothesis $P_\Obs\delta_\Source(Ju)\in\Hilb_\Obs$. In our setup this is
automatic: $\delta_\Source(Ju)\in\Hilb_\Source$ by definition of
$\Dom(\delta_\Source)$, and $P_\Obs=\E[\cdot\mid\mathcal F_\Obs]$ is an
$L^{2}$-contraction, so the conclusion always lies in $\Hilb_\Obs$.
\end{remark}

\begin{remark}[Reduction to exact intertwining]\label{rem:exact-special}
If $\delta_\Source(Ju)$ is $\mathcal F_\Obs$-measurable --- equivalently,
fixed by $P_\Obs$ --- then \cref{eq:proj-div} reduces to the exact
intertwining
\[
  \iota\,\delta_\Obs u \;=\; \delta_\Source(Ju),
\]
recovering the special case stated below in
\cref{cor:exact-intertwining}. For general lifts the divergence carries
a $P_\Obs$, which is the content of the present construction.
\end{remark}

\subsection*{Exact intertwining as a special case}

\begin{definition}[Exact divergence intertwining]\label{def:dintertwine}
The lift $J$ \emph{intertwines the divergences exactly} if there is a
subspace $\mathcal C\subseteq\Dom(\delta_\Obs)$, dense in $\Energy_\Obs$,
such that for every $u\in\mathcal C$, $Ju\in\Dom(\delta_\Source)$,
$\delta_\Source(Ju)$ is $\mathcal F_\Obs$-measurable, and
\[
  \iota\,\delta_\Obs u \;=\; \delta_\Source(Ju).
\]
\end{definition}

\begin{corollary}[Exact divergence intertwining]\label{cor:exact-intertwining}
Suppose $D_\Obs$ is the compression derivative of
\cref{thm:compression} and the inclusion $J$ intertwines the
divergences exactly in the sense of \cref{def:dintertwine}, with dense
subspace $\mathcal C\subseteq\Dom(\delta_\Obs)$. Then on $\mathcal C$
the projected-divergence formula \eqref{eq:proj-div} of
\cref{thm:proj-div} reduces to the exact identity
\[
  \iota\,\delta_\Obs u \;=\; \delta_\Source(Ju),\qquad u\in\mathcal C,
\]
without the conditional projection $P_\Obs$.
\end{corollary}

\begin{proof}
By \cref{thm:proj-div}, $\delta_\Obs u = P_\Obs\delta_\Source(Ju)$ for
$u\in\mathcal C$. Exact intertwining (\cref{def:dintertwine}) supplies
the $\mathcal F_\Obs$-measurability of $\delta_\Source(Ju)$, so
$P_\Obs\delta_\Source(Ju)=\delta_\Source(Ju)$ as elements of
$\Hilb_\Source$. Applying the inclusion $\iota$ to the identity
$\delta_\Obs u=P_\Obs\delta_\Source(Ju)$ in $\Hilb_\Obs$ gives the
displayed equation in $\Hilb_\Source$.
\end{proof}

\subsection*{Domain caveats}

\begin{remark}[Density of $\mathcal C_\Obs$ is non-trivial]
\label{rem:density-caveat}
Without \cref{ass:dense} the compression still produces a closable
operator on a non-dense domain, but $\delta_\Obs=D_\Obs^{*}$ is no longer
guaranteed to be densely defined and the duality machinery is lost. For
abstract Wiener spaces with a smooth source datum, density is typically
inherited from smooth cylindrical functionals; for composite processes
of the type considered in~\cite{IBMPaper} it is more delicate and is
treated there as an analytic task.
\end{remark}

\begin{remark}[Graph characterisation]\label{rem:minimality}
$D_\Obs$ is the unique closed operator
$\Hilb_\Obs\to\Energy_\Obs$ whose graph is the closure of
$\{(F,\,D_\Source(\iota F)):F\in\mathcal C_\Obs\}$ in
$\Hilb_\Obs\times\Energy_\Obs$. Choosing a different source
derivative produces, in general, a different observable derivative;
the compression is a functor from source-data to observable-data, not
a canonical construction independent of the source.
\end{remark}

\begin{remark}[Energy norm transfer]\label{rem:energy-norm}
Since $J$ is an isometric inclusion, for every $F\in\Dom(D_\Obs)$,
\[
  \|D_\Obs F\|_{\Energy_\Obs}
  \;=\; \|D_\Source(\iota F)\|_{\Energy_\Source}.
\]
The intrinsic energy is strictly the source energy restricted to those
modes that arise from observable functionals --- not a coarsened or
contracted version of it.
\end{remark}

\subsection*{Higher-order compression}

\begin{definition}[Second-order observable smooth core]
\label{def:obs-core-2}
Suppose a closed densely defined second-order source derivative
\[
  D_\Source^{(2)}\colon\Dom(D_\Source^{(2)})\subseteq\Hilb_\Source
   \longrightarrow\Energy_\Source^{(2)}
\]
into a second-order source energy $\Energy_\Source^{(2)}$ is given (in
the Gaussian setting, $D_\Source^{(2)}=D_\Source\otimes D_\Source$ on
the second source Sobolev space, with the appropriate tensor target).
The \emph{second-order observable smooth core} is
\[
  \mathcal C_\Obs^{(2)}\coloneqq
   \bigl\{F\in\Hilb_\Obs:\iota F\in\Dom(D_\Source^{(2)})\bigr\}.
\]
\end{definition}

\begin{definition}[Candidate second-order compression operator and energy]
\label{def:compression-2}
For $F\in\mathcal C_\Obs^{(2)}$, set
$D_\Obs^{(2),0}F\coloneqq D_\Source^{(2)}(\iota F)
\in\Energy_\Source^{(2)}$.
The \emph{candidate second-order observable energy} is
\[
  \Energy_\Obs^{(2)}
  \;\coloneqq\;
  \overline{\bigl\{\,D_\Source^{(2)}(\iota F):F\in\mathcal C_\Obs^{(2)}\,\bigr\}}
   ^{\,\Energy_\Source^{(2)}}
  \;\subseteq\;\Energy_\Source^{(2)},
\]
with isometric inclusion
$J^{(2)}\colon\Energy_\Obs^{(2)}\hookrightarrow\Energy_\Source^{(2)}$.
\end{definition}

\begin{proposition}[Second-order observable compression]
\label{prop:compression-2}
Assume $D_\Source^{(2)}$ is closed and densely defined and that
$\mathcal C_\Obs^{(2)}$ is dense in $\Hilb_\Obs$. Then
$D_\Obs^{(2),0}\colon\mathcal C_\Obs^{(2)}\subseteq\Hilb_\Obs
\to\Energy_\Obs^{(2)}$ is closed on the maximal pullback domain
$\mathcal C_\Obs^{(2)}$; we denote it $D_\Obs^{(2)}$. Any restriction
to a subcore $\mathcal C'\subseteq\mathcal C_\Obs^{(2)}$ dense in
$\Hilb_\Obs$ is closable, and the closure is contained in
$D_\Obs^{(2)}$.
\end{proposition}

\begin{proof}
Identical structurally to \cref{thm:compression}: if
$F_n\in\mathcal C_\Obs^{(2)}$ with $F_n\to F$ in $\Hilb_\Obs$ and
$D_\Obs^{(2),0}F_n\to h$ in $\Energy_\Obs^{(2)}$, the isometric
inclusions $\iota$ and $J^{(2)}$ give $\iota F_n\to\iota F$ in
$\Hilb_\Source$ and $J^{(2)}D_\Obs^{(2),0}F_n=D_\Source^{(2)}(\iota F_n)
\to J^{(2)}h$ in $\Energy_\Source^{(2)}$. Closedness of
$D_\Source^{(2)}$ yields $\iota F\in\Dom(D_\Source^{(2)})$, i.e.,
$F\in\mathcal C_\Obs^{(2)}$, and $D_\Obs^{(2),0}F=h$. Closability of
subcore restrictions is the same
argument applied to sequences in $\mathcal C'$.
\end{proof}

\Cref{prop:compression-2} produces a closed operator; it does not
produce a generator. Identifying the action of $D_\Obs^{(2)}$ on a
useful test class and converting it into an It\^o-type correction is
the next analytic problem, structured by the energy-order
classification of \cref{sec:energy-order}.

\section{The conditional operator-It\^o transfer principle}\label{sec:transfer}

We layer the framework one step further. Combining the first-order
observable divergence $\delta_\Obs$ (\cref{thm:proj-div}) with a
source-level It\^o expansion, we state and prove an abstract
\emph{conditional} transfer theorem showing how a source-level It\^o
expansion descends to an observable-level one.

Let $(\Hilb_\Source,\Energy_\Source,D_\Source)$ and
$(\Hilb_\Obs,\Energy_\Obs,D_\Obs)$ be the source and observable energy
data of \cref{sec:compat}, with isometric inclusions $\iota$ and $J$
and observable projection
$P_\Obs=\iota^*=\E[\,\cdot\mid\mathcal F_\Obs]$.
Suppose an observable process $U=(U_t)_{t\ge 0}$ and a smooth class of
test functions admit a \emph{source-level} It\^o-type expansion
\begin{equation}\label{eq:source-ito}
  f(U_t)-f(U_0)\;=\;\delta_\Source\bigl(A_f(t)\bigr)\;+\;R_f(t),
  \qquad A_f(t)\in\Dom(\delta_\Source),
\end{equation}
on the joint source $\sigma$-algebra $\mathcal F_\Source$, with
remainder $R_f(t)\in L^{2}(\mathcal F_\Source)$.

\begin{theorem}[Conditional operator-It\^o transfer]\label{thm:transfer}
Assume \eqref{eq:source-ito} holds with $A_f(t)\in J(\Energy_\Obs)$,
so that there exists $B_f(t)\in\Energy_\Obs$ with $JB_f(t)=A_f(t)$; in
particular $B_f(t)$ lies in the lifted source-divergence domain
$\{u\in\Energy_\Obs:Ju\in\Dom(\delta_\Source)\}$, so
\cref{thm:proj-div} applies. Suppose further that $f(U_t)-f(U_0)$ is
$\Prob$-a.s.\ $\mathcal F_\Obs$-measurable. Then
\begin{equation}\label{eq:obs-ito}
  f(U_t)-f(U_0)\;=\;\delta_\Obs\bigl(B_f(t)\bigr)\;+\;P_\Obs R_f(t),
\end{equation}
with $\delta_\Obs(B_f(t))=P_\Obs\,\delta_\Source(A_f(t))$.
\end{theorem}

\begin{proof}
By \cref{thm:proj-div} applied to $B_f(t)$,
$\delta_\Obs(B_f(t))=P_\Obs\,\delta_\Source(JB_f(t))
=P_\Obs\,\delta_\Source(A_f(t))$. Apply $P_\Obs$ to both sides of
\eqref{eq:source-ito}: since $f(U_t)-f(U_0)$ is
$\mathcal F_\Obs$-measurable, $P_\Obs[f(U_t)-f(U_0)]=f(U_t)-f(U_0)$;
linearity of $P_\Obs$ on the right gives
$\delta_\Obs(B_f(t))+P_\Obs R_f(t)$.
\end{proof}

\Cref{thm:transfer} does \emph{not} prove an It\^o formula for $U$ on
its own. It converts a source-level identity \eqref{eq:source-ito}
into an observable-level identity \eqref{eq:obs-ito}, conditional on
the source divergence integrand falling in the lifted observable image
$J(\Energy_\Obs)$. The practical content is a separation of duties:
the source side supplies the It\^o expansion; the compression
construction (\cref{thm:compression,thm:proj-div}) supplies the
observable lift; the correction term is the observable projection of
the source remainder.

The theorem does not construct $A_f(t)$, $B_f(t)$, or $R_f(t)$ in a
given example. The source-level It\^o expansion \eqref{eq:source-ito}
is the input to the transfer; producing it for a specific composite
observable is part of the application paper. The transfer's content is
that, given the source-level identity and the lifting hypothesis
$A_f(t)\in J(\Energy_\Obs)$, the conditional-expectation formula of
\cref{thm:proj-div} yields the observable-level identity automatically,
with the observable correction equal to the conditional expectation of
the source remainder.

Two hypotheses of \cref{thm:transfer} are non-trivial in concrete
examples:
\begin{itemize}
\item[(T1)] $A_f(t)\in J(\Energy_\Obs)$ --- the source divergence
  integrand lifts through the observable inclusion;
\item[(T2)] $f(U_t)-f(U_0)$ is $\mathcal F_\Obs$-measurable.
\end{itemize}
(T2) is automatic when $U$ is itself observable in $\mathcal F_\Obs$,
which is the natural setting. (T1) is a genuine compatibility
requirement on the source It\^o expansion: it asserts that the
source-level martingale-direction integrand $A_f(t)$, while a priori
an element of $\Energy_\Source$, in fact arises by lifting an
observable element $B_f(t)\in\Energy_\Obs$. When this fails, the
source-level expansion \eqref{eq:source-ito} does not descend to a
strictly observable identity; what one can still extract is a
partially-projected identity in $L^{2}(\mathcal F_\Source)$, which is
not what the framework targets.

The transfer principle is first-order: it converts $A_f(t)$ in
$\Energy_\Source$ into $B_f(t)$ in $\Energy_\Obs$, and the
remainder $R_f(t)$ in $L^{2}(\mathcal F_\Source)$ into
$P_\Obs R_f(t)$ in $L^{2}(\mathcal F_\Obs)$. The structural shape of
the remainder --- in particular the observable correction order $q$
at which it should be read --- is not part of the transfer; that is
the content of the energy-order classification of
\cref{sec:energy-order}. In the quadratic-energy case
($q=2$, operator rank $r=1$),
$R_f(t)=\tfrac12\int_0^t f''(U_s)d\langle U\rangle_s$ recovers the
semimartingale It\^o correction; in higher-order cases ($q\ge 4$,
operator rank $r\ge 2$), the remainder is naturally read against a
higher-order observable correction energy $\Energy_\Obs^{[q]}$
generated by an operator-rank-$r$ datum $D_\Source^{(r)}$ obtained by
compressing a higher-rank source derivative in the sense of
\cref{def:compression-2,prop:compression-2}. The remaining structural
content of the framework is the principled classification, by
\emph{observable correction order}, of the candidate identities
produced by the
transfer, to which we turn next.

\section{Energy-order classification}\label{sec:energy-order}

The central claim of the energy-first reorganization is that It\^o's
formula is a parametric family indexed by \emph{observable correction
order}. The parametric It\^o identity
\eqref{eq:abstract-ito-hierarchy} from the abstract,
\begin{equation*}
  f(U_t)-f(U_0)\;=\;\delta\bigl(B_f(t)\bigr)\;+\;\mathrm{Corr}_f^{[q]}(t),
  \qquad\mathrm{Corr}_f^{[q]}(t)\in\bigl(\Energy^{[q]}\bigr)^{*},
\end{equation*}
asserts that, under the source-expansion and lift hypotheses of the
transfer theorem (\cref{thm:transfer}), the transfer principle
produces an It\^o-type expansion of the observable whose centred part
is a divergence in the observable energy and whose correction lives in
the dual of an order-$q$ observable correction energy space, with the
integer $q$ determined by the noise composition (for the $q=2$ vs
$q\ge 4$ realisation of $(\Energy^{[q]})^{*}$, see
\cref{prop:quartic-dual}). The order $q$ is proposed as the natural organising invariant of the
identity, in place of semimartingale quadratic variation: the
quadratic-variation case is recovered at $q=2$, and higher-order
composite-noise observables produce instances at $q\ge 4$. This section develops the operator-theoretic content of the
parametric identity, recovers the classical case at $q=2$ (operator
rank $r=1$), and records the higher-order witnesses at $q\ge 4$.

\begin{remark}[Notation convention: $q$ versus $r$]
\label{rem:q-r-convention}
Throughout this paper, two distinct integers index a Level-3
It\^o-type identity:
\begin{itemize}
\item[$q$ ---] the \emph{observable correction order}: the order of
  the highest test-function derivative $f^{(q)}$ appearing in the
  correction $\mathrm{Corr}_f^{[q]}$, equivalently the spatial order
  of the PDE shadow of $X$ in the Level-1 observable It\^o identity
  \eqref{eq:obs-ito-identity}. Bracket superscripts $[q]$ index
  correction-order objects: $\mathrm{Corr}_f^{[q]}$, $\Energy^{[q]}$,
  $C_f^{[q]}$.
\item[$r$ ---] the \emph{operator rank}: the rank of the
  operator-derivative datum $D^{(r)}\colon\Dom(D^{(r)})\to
  \Energy^{(r)}$ whose carr\'e form generates the correction.
  Parenthesis superscripts $(r)$ index operator-rank objects:
  $D^{(r)}$, $\Energy^{(r)}$, $\Carre^{(r)}$. (The term ``rank'' is
  used here in the order-of-derivative sense and does not refer to the
  dimension of the range of $D^{(r)}$.)
\end{itemize}
In the even-order tensor cases considered in this paper,
\begin{equation*}
  q\;=\;2r.
\end{equation*}
Explicitly: classical semimartingale It\^o is the case $q=2$, $r=1$
(a single operator derivative $D$, ordinary carr\'e
$\Carre(X,X)=\langle DX,DX\rangle_\Energy$, second-order PDE shadow);
iterated Brownian motion is the case $q=4$, $r=2$ (a rank-$2$
operator derivative $D^{(2)}$, rank-$2$ carr\'e built from $D^{(2)}$,
fourth-order PDE shadow). The two indices are not interchangeable: the
classical correction $\tfrac12 f''(X_s)\,d\langle X\rangle_s$ is built
from a \emph{first}-rank operator derivative even though its PDE
shadow is of order two.
\end{remark}

To anchor the development, recall that the classical semimartingale
It\^o formula reads
\begin{equation}\label{eq:classical-ito}
  f(X_t)-f(X_0)
  \;=\;\int_0^t f'(X_s)\,dX_s
   \;+\;\tfrac12\!\int_0^t f''(X_s)\,d\langle X\rangle_s,
\end{equation}
and the structural content of \eqref{eq:classical-ito} is that the
correction is encoded by the quadratic-variation bracket
$d\langle X\rangle$. This bracket is the trace of a
\emph{second-order} object: in the bracket-compatible energy
representations covered by \cref{thm:k2-recovery},
$d\langle X\rangle$ corresponds to the diagonal of the carr\'e du
champ $\Carre(F,F)=\langle DF,DF\rangle_\Energy$ for $F=X_t$,
i.e.~the trace of $DX\otimes DX$ against the source pairing. The
energy-first framework reads the same identity in terms of the order
of the energy form rather than of the quadratic-variation bracket.

\begin{definition}[Observable correction order and operator rank]
\label{def:energy-order}
Suppose an observable process $X$ admits a Level-3 lift in the sense
of \cref{eq:operator-ito}:
\begin{equation}\label{eq:order-q-ito}
  f(X_t)-f(X_0)\;=\;\delta(U_f(X_\bullet))+C_f^{[q]}(X_\bullet).
\end{equation}
The lift has \emph{observable correction order $q$} if the highest
test-function derivative entering the correction $C_f^{[q]}$ is
$f^{(q)}$ on the chosen test class --- equivalently, if the leading
differential part of the PDE shadow $\mathcal G_t$ of $X$ (the
differential operator appearing in the Level-1 observable It\^o
identity~\eqref{eq:obs-ito-identity}) is of order $q$
(boundary/local-time contributions are absorbed into the lower-order
terms of \cref{rem:schematic-correction}).
It has \emph{operator rank $r$} if the correction takes the schematic
form~\eqref{eq:schematic-correction} for an $r$-fold
operator-derivative datum
$D^{(r)}\colon\Dom(D^{(r)})\subseteq\Hilb\to\Energy^{(r)}$ with
carr\'e form
\[
  \Carre^{(r)}(X,\dots,X)
   \;=\;\langle D^{(r)}X,\,D^{(r)}X\rangle_{\Energy^{(r)}}.
\]
In the even-order tensor cases considered in this paper, the two
indices are related by $q=2r$, which determines $r$ uniquely from $q$;
we take this as the convention throughout. The reference cases of
this classification are
\begin{itemize}
\item the classical semimartingale It\^o formula: $q=2$, $r=1$,
  generated by a single operator derivative $D$ with ordinary carr\'e
  $\Carre(X,X)=\langle DX,DX\rangle_\Energy$;
\item iterated Brownian motion: $q=4$, $r=2$, generated by a rank-$2$
  operator derivative $D^{(2)}$ with rank-$2$ carr\'e
  $\Carre^{(2)}(X,X)=\langle D^{(2)}X,D^{(2)}X\rangle_{\Energy^{(2)}}$.
\end{itemize}
See \cref{rem:q-r-convention} for the firm notational separation of
$q$ (bracket superscripts) from $r$ (parenthesis superscripts).
\end{definition}

\begin{remark}[Schematic form of the leading correction]
\label{rem:schematic-correction}
In the basic scalar even-order tensor examples treated in this paper,
the leading correction $C_f^{[q]}$ has the schematic form
\begin{equation}\label{eq:schematic-correction}
  C_f^{[q]}(X_\bullet)
   \;\sim\;\int_0^t f^{(q)}(X_s)\,
        \Carre^{(r)}(X,\dots,X)_{s}\,ds
   \;+\;\text{(boundary / local-time / lower-order terms)},
\end{equation}
with $q=2r$, but the numerical coefficient (e.g.\ $\tfrac12$ in the
classical $q=2$ case, $\tfrac18$ in the IBM $q=4$ case) and the
specific boundary or local-time contributions are model-dependent and
are read off in each instance from the source-level It\^o expansion;
\cref{def:energy-order} does not commit to a universal coefficient
formula.
\end{remark}

\begin{program}[Observable correction-order classification]\label{prog:energy-order}
The natural organising invariant of an It\^o-type identity is the
\emph{observable correction order $q$} that produces the correction
$C_f^{[q]}$ (equivalently, the order of the PDE shadow of $X$), not
the existence of a quadratic-variation bracket. The classical
semimartingale It\^o formula is the $q=2$ instance of this
classification (operator rank $r=1$); observables whose pathwise
correction is determined by a higher-rank operator-derivative datum
have observable correction order $q>2$, and the PDE shadow of the
correction is a differential operator of order $q$.
\end{program}

The thesis is a programme rather than a theorem because, while
\cref{def:energy-order} makes the correction order well-defined once
an energy datum and a Level-3 lift are chosen, identifying the correct
energy datum and verifying the lift on a useful test class are
problem-specific analytic tasks. We record three structural
consequences: a recovery theorem in the classical $q=2$, $r=1$ case
(treated below), an existence statement for higher-order witnesses,
and a target-topology theorem for non-classical corrections (both
treated thereafter).

\subsection*{The classical case: \texorpdfstring{$q=2$, $r=1$}{q=2, r=1} recovers semimartingale It\^o--Tanaka}

At observable correction order $q=2$ with operator rank $r=1$, the
parametric identity \eqref{eq:abstract-ito-hierarchy} specialises to
the classical semimartingale It\^o formula. The order-$2$ observable
correction energy $\Energy^{[2]}$ is the dual of the ordinary
carr\'e form
\[
  \Carre(X,X)\;=\;\langle DX,DX\rangle_{\Energy},
\]
which bracket compatibility identifies with $d\langle X\rangle$. The
correction $\mathrm{Corr}_f^{[2]}(t)$ in
\eqref{eq:abstract-ito-hierarchy} reads
\[
  \mathrm{Corr}_f^{[2]}(t)
   \;=\;\tfrac12\!\int_0^{t}f''(X_s)\,d\langle X\rangle_s,
\]
generated by a single operator derivative $D$ (operator rank $r=1$).
The quadratic-variation bracket is therefore not absent in the
energy-first framework; it is the carr\'e form of the rank-$1$
operator datum, and the $q=2$ instance of the parametric identity
coincides with classical It\^o--Tanaka.

In the following, $D$ is not assumed to be the standard Malliavin
derivative; it is the operator derivative associated with the chosen
energy datum $(\Hilb,\Energy,D)$. The compatibility hypotheses below
are precisely what identify its local carr\'e with the semimartingale
bracket.

\begin{theorem}[Classical semimartingale It\^o as the $q=2$, $r=1$
instance]\label{thm:k2-recovery}
Let $X$ be a continuous semimartingale on a complete filtered
probability space $(\Omega,\mathcal F_t,\Prob)$, with predictable
quadratic variation $\langle X\rangle_t$. Let $(\Hilb,\Energy,D)$ be a
bracket-compatible operator-energy datum on which $X_t\in\Dom(D)$ for
every $t\ge 0$ and the chain rule $D(f(X_t))=f'(X_t)DX_t$ holds for
$f\in C^{2}(\R)$. Assume $s\mapsto DX_s$ is a predictable
$\Energy$-valued process with
$\int_{0}^{t}\|DX_s\|_{\Energy}^{2}\,ds<\infty$ a.s., and define
the predictable random measure
\[
  d\Carre^{D}_{X}(\omega,ds)\;\coloneqq\;\langle DX_s,DX_s\rangle_{\Energy}\,ds.
\]
Assume further the following two compatibility hypotheses on the
class of test integrands $f'(X_\bullet)DX$ under consideration:
\begin{itemize}
\item[(B)] \emph{Bracket compatibility (measure form):}
  \begin{equation}\label{eq:bracket-compat}
    d\Carre^{D}_{X}(\omega,ds) \;=\; d\langle X\rangle_s(\omega)
  \end{equation}
  as predictable random measures on $\R_{+}$. (If $\langle X\rangle$
  is absolutely continuous in $t$ with density
  $\tfrac{d\langle X\rangle_s}{ds}$, equivalently
  $\langle DX_s,DX_s\rangle_{\Energy}=\tfrac{d\langle X\rangle_s}{ds}$
  a.e.~$s$.)
\item[(D)] \emph{Divergence compatibility:} $f'(X_\bullet)DX
  \in\Dom(\delta)$ for $f\in C^{2}(\R)$ in the test class, and
  \begin{equation}\label{eq:div-compat}
    \delta\bigl(f'(X_\bullet)DX\bigr)
     \;=\;\int_0^t f'(X_s)\,dX_s.
  \end{equation}
\end{itemize}
Then the parametric It\^o identity \eqref{eq:abstract-ito-hierarchy}
at $q=2$ (operator rank $r=1$) recovers the classical semimartingale
It\^o formula
\begin{equation*}
  f(X_t)-f(X_0) \;=\; \int_0^t f'(X_s)\,dX_s
   \;+\;\tfrac12\!\int_0^t f''(X_s)\,d\langle X\rangle_s,
\end{equation*}
and $X$ has observable correction order $q=2$ and operator rank $r=1$
in the sense of \cref{def:energy-order}, with
$C_f^{[2]}=\tfrac12\int_0^t f''(X_s)\,d\langle X\rangle_s$.
\end{theorem}

\begin{proof}
By the classical semimartingale It\^o formula for continuous
semimartingales,
\[
  f(X_t)-f(X_0)
   \;=\;\int_0^{t}f'(X_s)\,dX_s
   \;+\;\tfrac12\!\int_0^{t}f''(X_s)\,d\langle X\rangle_s.
\]
By divergence-compatibility~(D), the first term equals
$\delta(f'(X_\bullet)DX)$, so it has the form $\delta(B_f(t))$ with
$B_f(t)=f'(X_\bullet)DX\in\Energy$.
By bracket-compatibility~(B), $d\langle X\rangle_s=d\Carre^{D}_{X}(\omega,ds)$
as predictable random measures, so the second term equals
\[
  \tfrac12\!\int_0^{t}f''(X_s)\,d\Carre^{D}_{X}(\omega,ds)
   \;=\;\mathrm{Corr}_f^{[2]}(t),
\]
paired against $f''$ in $(\Energy^{[2]})^{*}$. The $f''$-pairing
fixes $q=2$; the single operator derivative $D$ fixes $r=1$.
Hypothesis~(D) is the Skorohod/It\^o agreement on the chosen
predictable test class; on classical Wiener space (the canonical
Brownian/Malliavin specialisation, with $D$ the Malliavin derivative
and $\Energy=L^{2}(\R_{+})$) this is the standard fact recorded
in~\cite{Nualart}.
\end{proof}

\Cref{thm:k2-recovery} clarifies the relationship between
the classical construction and the energy-first reorganization. The
classical organising principle for It\^o calculus is the existence
of a quadratic-variation bracket $\langle X\rangle$. When $X$ is not
a semimartingale in its natural filtration $\sigma(X_s:s\le t)$, it
does not carry the finite predictable bracket structure required by
the classical semimartingale It\^o formula, and the classical theory
does not produce an It\^o formula at that level. The energy-order
classification refuses to take the absence of an admissible
$\langle X\rangle$ as obstructive: a process can have observable
correction order $q\ne 2$, in which case its pathwise correction is
built from a higher-rank operator-derivative carr\'e form
$\Carre^{(r)}$ (with $q=2r$ in the even tensor case) rather than from
a quadratic bracket. The non-semimartingale obstruction is
reinterpreted as a failure of the $q=2$ coordinate, rather than as an
obstruction to all It\^o-type structure.
The energy-first reorganization proposes the observable correction
order $q$ as the natural invariant of an It\^o-type identity, replacing
semimartingale quadratic variation. The classical theory is the case
$q=2$, $r=1$; the framework provides the operator-theoretic
infrastructure for higher-order cases, with the IBM case ($q=4$,
$r=2$) developed conditionally in the companion paper
(\cref{cor:k4-witness}).

\subsection*{Higher-order cases: composite processes and the
witnesses at \texorpdfstring{$q\ge 4$}{q>=4}}

Beyond $q=2$, the parametric identity
\eqref{eq:abstract-ito-hierarchy} predicts that processes whose
pathwise corrections fail to live in $L^{2}(\Omega)$ at the
second-order level should be readable at higher correction order. In
a concrete example, the correction order of a process is read off the
Level-1 observable It\^o identity: \cref{def:obs-ito-identity}
records the observable correction operator $\mathcal G_t$, and the
candidate observable correction order $q$ equals the order of
$\mathcal G_t$ as a differential operator on the test class. The
associated operator rank is $r=q/2$ in the even-order tensor cases
considered here (see \cref{rem:q-r-convention}). The diagnostic is
therefore PDE-shadow data: one differentiates $\E[f(X_t)]$ in $t$ on a
test class and reads the order of the resulting operator. The
energy-first reorganization then asks what operator-rank-$r$ datum and
what order-$q$ Level-3 lift carry that PDE shadow back to the pathwise
level. The next two results record a concrete $q=4$ witness conditional
on the companion paper~\cite{IBMPaper}, and the abstract
target-topology structure required when the raw $L^{2}(\Omega)$
correction fails.

\begin{corollary}[$q=4$, $r=2$ witness, conditional
on~\cite{IBMPaper}]\label{cor:k4-witness}
Conditional on the mollified It\^o formula proved
in~\cite[Theorem~7.1]{IBMPaper}, iterated Brownian motion
$\IBM_t=X_{|Y_t|}$ is a $q=4$, $r=2$ witness of the energy-order
classification of \cref{sec:energy-order}: the observable correction
energy $\Energy_\IBM^{[4]}$ is fourth-order, and the parametric It\^o
identity~\eqref{eq:abstract-ito-hierarchy} is realised at $q=4$ with
the specific correction $\mathrm{Corr}_f^{[4]}(t)$ computed
in~\cite[Theorem~7.1]{IBMPaper}, generated by a rank-$2$ operator
derivative $D_\IBM^{(2)}$ (operator rank $r=2$). The associated
observable correction operator is $\mathcal G_t=\tfrac18\partial_x^{4}$
plus a boundary term, so by the diagnostic above the observable
correction order is $q=4$, consistent with $q=2r$. Here $q$ is read
off the leading differential part of $\mathcal G_t$; the boundary
contribution is absorbed into the lower-order/local-time terms of the
schematic correction (\cref{rem:schematic-correction}). In particular IBM
lies strictly outside the bracket-compatible class of
\cref{thm:k2-recovery}: it is not a semimartingale in its natural
filtration $\sigma(\IBM_s:s\le t)$ and does not carry the finite
predictable bracket structure required by the classical semimartingale
It\^o formula.
\end{corollary}

\begin{proof}
Direct application of~\cite[Theorem~7.1]{IBMPaper} together with the
energy-order classification of \cref{def:energy-order}: the cited
theorem supplies the source-level It\^o expansion of $f(\IBM_t)$
generated by an operator-rank-$2$ datum, and \cref{thm:transfer}
converts it into the observable identity whose correction is read at
$q=4$ in $(\Energy_\IBM^{[4]})^{*}$.
\end{proof}

\begin{proposition}[Rigged hierarchy via Gelfand triple]
\label{prop:quartic-dual}
Assume a dense, continuously embedded subspace
$\Energy^{[q],+}\hookrightarrow\Energy^{[q]}$.
Then the standard Gelfand-triple
construction~\cite{Berezansky,HidaKuoPotthoff} produces the rigged
hierarchy
\begin{equation}\label{eq:rigged-4-abstract}
  \Energy^{[q],+}\;\subset\;\Energy^{[q]}\;\subset\;\Energy^{[q],-},
\end{equation}
where $\Energy^{[q],-}$ is the antidual of $\Energy^{[q],+}$ in the
topology induced by the embedding. When the raw
$L^{2}(\Omega)$-correction of an order-$q$ Level-3 lift does not exist,
a natural candidate target topology for the pathwise correction is
$\Energy^{[q],-}$, paired against test integrands from
$\Energy^{[q],+}$.
\end{proposition}

\begin{proof}
The construction of~\eqref{eq:rigged-4-abstract} is the standard
Gelfand-triple construction; see~\cite[Ch.~1]{Berezansky}
or~\cite[Sec.~3]{HidaKuoPotthoff}. Concrete realisations of
$\Energy^{[q],+}$ (Hida-distributional, weighted-Sobolev, or
polynomial-chaos completions) are constructed in the companion papers:
for iterated Brownian motion the quartic case $q=4$ uses
$\Energy_\IBM^{[4]}=\Energy_\IBM^{(2)}\otimes\Energy_\IBM^{(2)}$
(correction energy at $q=4$ realised as the tensor square of the
rank-$r=2$ operator energy $\Energy_\IBM^{(2)}$) with the rigged
hierarchy of~\cite[\S5]{IBMPaper}; the abstract framework requires
only the existence of such a triple.
\end{proof}

The rigged hierarchy
$\Energy^{[q],+}\subset\Energy^{[q]}\subset\Energy^{[q],-}$ used
here is an instance of the standard Gelfand-triple construction: a
dense, continuously embedded subspace $\Energy^{[q],+}$ supplies test
integrands, and the antidual $\Energy^{[q],-}$ supplies the topology
in which higher-order corrections naturally live when they fail to be
square-integrable. Concrete realisations of the triple use
Hida-distributional, weighted-Sobolev, or polynomial-chaos
completions depending on the noise structure; in the iterated
Brownian motion case the rigged hierarchy of~\cite[\S5]{IBMPaper}
uses weighted-Sobolev spaces $H^{s}$ in the spatial scenery direction
combined with chaos completions in the clock direction. The framework
paper requires only the existence of such a triple, with continuous
embedding and closability of the divergence on the test subspace; the
construction itself is analytical and is treated case by case.

The natural higher-order witnesses of the classification are
composite processes. The worked example of~\cite{IBMPaper} exhibits
a fourth-order PDE shadow on the composite process
$\IBM_t=X_{|Y_t|}$: by \cref{cor:k4-witness} and
\cref{prop:quartic-dual}, $\IBM$ is the $q=4$ instance of the
parametric identity \eqref{eq:abstract-ito-hierarchy} --- PDE-shadow
order four, with the order-$q$ correction
$\mathrm{Corr}_f^{[4]}(t)$ living in the dual
$(\Energy_\IBM^{[4]})^{*}$ of the quartic composite observable
correction energy
$\Energy_\IBM^{[4]}=\Energy_\IBM^{(2)}\otimes\Energy_\IBM^{(2)}$
rather than in $L^{2}(\Omega)$. The composite structure is what
produces the fourth-order signature: the clock It\^o formula behind
the IBM identity is second-order in the Brownian-time variable, but
because its coefficients are derivatives of Brownian scenery the
observable correction has fourth-order PDE shadow. The mechanism, in
the q/r language of \cref{rem:q-r-convention}, is a single one: the
two rank-$1$ source Malliavin calculi $(D_X,\delta_X)$ and
$(D_Y,\delta_Y)$ compose into a rank-$2$ observable operator-derivative
datum $D_\IBM^{(2)}$, whose carr\'e form $\Carre^{(2)}_\IBM$ produces
the $q=4$ correction energy
$\Energy_\IBM^{[4]}=\Energy_\IBM^{(2)}\otimes\Energy_\IBM^{(2)}$ and
the fourth-order PDE shadow. We expect a wider family of composite, subordinated,
time-changed, and iterated processes to populate the classification
at orders $q\ge 4$, with the framework of
\cref{thm:compression,thm:proj-div,thm:transfer} producing the
operator/energy objects and the parametric identity
\eqref{eq:abstract-ito-hierarchy} reading off the order-$q$
correction in each such case.

\section{Outlook: the worked example in the companion paper}\label{sec:outlook}

The framework developed in
\cref{thm:compression,thm:proj-div,thm:transfer} is, on its own, an
abstract recipe: it gives a compression construction, a projected
divergence formula, and a conditional operator-It\^o transfer
principle. To exhibit its reach in a regime where classical methods
do not apply, the companion paper~\cite{IBMPaper} carries out the
construction in detail for the canonical non-semimartingale
composite-noise example: \emph{iterated Brownian motion}
$\IBM_t=X_{|Y_t|}$, with two-sided Brownian scenery $X$ and an
independent Brownian clock $Y$. Three features make IBM diagnostically
useful as a witness of the energy-first framework. Its PDE shadow is
the bi-Laplacian/local-time PDE signature identified by
Allouba--Zheng~\cite{AllouZheng}, not the heat equation: IBM is not
Markov in its own filtration and its observable correction operator
is $\mathcal G_t=\tfrac18\partial_x^{4}$ plus a boundary term, so by
the energy-order classification of \cref{prog:energy-order} it is a
fourth-order observable It\^o process (PDE-shadow order four,
observable correction order $q=4$ and operator rank $r=2$ in the sense
of \cref{rem:q-r-convention}: a rank-$2$ operator derivative whose
carr\'e form, contracted with the clock, produces the fourth-order
spatial shadow, consistent with $q=2r$). It is also naturally a
\emph{composite} process: the source Malliavin calculi
$(D_X,\delta_X)$ and $(D_Y,\delta_Y)$ live on the two independent
Gaussian factors and provide an explicit atlas indexed by the chaos
bidegree $(m,n)$, as developed in~\cite{Fontes-IBMchaos} --- the
source datum is concrete, and the observable datum is obtained by the
compression construction of \cref{thm:compression}. Finally, IBM does
not fall under the classical semimartingale It\^o theory in its own
filtration: it is not a semimartingale in its natural filtration
$\sigma(\IBM_s:s\le t)$ and does not carry the finite predictable
bracket structure required by the classical semimartingale It\^o
formula, so the energy-order classification, rather than the
$\langle\cdot\rangle$-bracket framework, is the natural organising
tool.

The companion paper~\cite{IBMPaper} applies the compression
construction \cref{thm:compression} and the projected divergence
formula \cref{thm:proj-div} to the IBM source $(X,Y)$ and the
observable $\sigma$-algebra
$\mathcal F_\IBM\subseteq\mathcal F_X\vee\mathcal F_Y$, obtaining a
closed first-order intrinsic IBM derivative $D_\IBM$ on the maximal
pullback core $\mathcal C_\IBM$. The abstract framework is realised
concretely on a mollified occupation core, with explicit kernels and
quantitative $\varepsilon\to 0$ convergence in the scenery-energy
direction. The main pathwise theorem of~\cite{IBMPaper} is a
\emph{mollified pathwise IBM It\^o formula} at every $\varepsilon>0$,
with explicit pathwise correction expressed through a clock-Skorohod
divergence and a Tanaka-style boundary local-time term, and its
fourth-order PDE-shadow consistency in expectation as
$\varepsilon\to 0$. Beyond the mollified level, the companion paper
proves the $L^{2}(\Omega)$ obstruction to the unmollified pathwise
correction: the centred first-scenery-chaos projection of the
candidate correction has variance that diverges as
$\varepsilon\to 0$, via an explicit Stein-cancellation argument that
isolates a residual $\rho_\varepsilon'$-kernel driven by the
Bouleau--Yor spatial quadratic variation of the Brownian clock local
time. The structural target topology for the unmollified correction
is then the dual $\Energy_\IBM^{[4],-}$ of a quartic composite
observable correction energy
$\Energy_\IBM^{[4]}=\Energy_\IBM^{(2)}\otimes\Energy_\IBM^{(2)}$
(correction order $q=4$ realised as the tensor square of the rank-$2$
operator energy $\Energy_\IBM^{(2)}$), in exact accord with
\cref{prop:quartic-dual}. Each of these results is
an instance of the abstract framework proved here applied in a
specific composite-noise regime: the abstract framework supplies the
operator/energy structures and the conditional operator-It\^o
transfer; the worked example supplies the concrete analytic content
(mollifier estimates, Stein cancellation, Bouleau--Yor input) needed
to land a proved theorem at the mollified level and to isolate the
obstruction at the unmollified level.

The framework is not specific to IBM. Fractional Brownian motion,
Volterra and ambit processes, stable L\'evy processes, stochastic
volatility models, and a wider family of composite, subordinated,
time-changed, and iterated processes are natural candidates for the
energy-order classification: each provides a source energy datum, a
sub-$\sigma$-algebra encoding the observable, and a PDE shadow whose
order indicates the observable correction order $q$ of the observable.
We do not develop these applications here. The framework provides the
operator/energy infrastructure; specific applications follow the
pattern set by~\cite{IBMPaper} of identifying the relevant source
datum, applying \cref{thm:compression,thm:proj-div,thm:transfer},
and reading off the appropriate observable correction order $q$ from
the PDE shadow.

\section*{Acknowledgments}

AI tools were used as interactive assistants for drafting,
reorganizing, and generating mathematical arguments, under the
author's direction. The author is responsible for all mathematical
content and any errors.

\begin{thebibliography}{99}

\bibitem{AllouZheng}
H.~Allouba and W.~Zheng,
\emph{Brownian-time processes: the PDE connection and the
half-derivative generator},
Ann. Probab. \textbf{29} (2001), no.~4, 1780--1795.

\bibitem{Berezansky}
Yu.~M.~Berezansky and Yu.~G.~Kondratiev,
\emph{Spectral Methods in Infinite-Dimensional Analysis. Vol.~1},
Mathematical Physics and Applied Mathematics, vol.~12/1,
Kluwer Academic Publishers, Dordrecht, 1995.

\bibitem{BogachevGaussian}
V.~I.~Bogachev,
\emph{Gaussian Measures},
Mathematical Surveys and Monographs, vol.~62,
American Mathematical Society, Providence, RI, 1998.

\bibitem{BouleauHirsch}
N.~Bouleau and F.~Hirsch,
\emph{Dirichlet Forms and Analysis on Wiener Space},
de Gruyter Studies in Mathematics, vol.~14,
Walter de Gruyter \& Co., Berlin, 1991.

\bibitem{Fontes-IBMchaos}
R.~Fontes,
\emph{Source-resolved chaos expansions for iterated Brownian motion},
companion paper, in preparation, 2026.

\bibitem{Fontes-IBMnumerical}
R.~Fontes,
\emph{A numerical variance budget for iterated Brownian motion},
companion paper, in preparation, 2026.

\bibitem{Fontes-OperatorDerivative}
R.~Fontes,
\emph{The operator derivative in continuous stochastic calculus:
a Hilbert energy space framework},
preprint, 2026.

\bibitem{Fontes-PDEgenerators}
R.~Fontes,
\emph{PDE generators and energy corrections for Brownian-time processes},
companion paper, in preparation, 2026.

\bibitem{FukushimaOshimaTakeda}
M.~Fukushima, Y.~Oshima, and M.~Takeda,
\emph{Dirichlet Forms and Symmetric Markov Processes},
2nd ed., de Gruyter Studies in Mathematics, vol.~19,
Walter de Gruyter \& Co., Berlin, 2011.

\bibitem{HidaKuoPotthoff}
T.~Hida, H.-H.~Kuo, J.~Potthoff, and L.~Streit,
\emph{White Noise: An Infinite Dimensional Calculus},
Mathematics and Its Applications, vol.~253,
Kluwer Academic Publishers, Dordrecht, 1993.

\bibitem{IBMPaper}
R.~Fontes,
\emph{Iterated Brownian Motion: An It\^o Formula in the Energy-First
Framework} (working title),
companion paper, in preparation, 2026. [Placeholder citation key to be
fixed at submission.]

\bibitem{StickyFlowNoiseLift}
R.~Fontes,
\emph{Sticky-Flow Noise Lifts in the Energy-First Framework: the
Le~Jan--Raimond Case} (working title),
companion paper, in preparation, 2026. [Placeholder citation key to be
fixed at submission.]

\bibitem{Janson}
S.~Janson,
\emph{Gaussian Hilbert Spaces},
Cambridge Tracts in Mathematics, vol.~129,
Cambridge University Press, Cambridge, 1997.

\bibitem{Nualart}
D.~Nualart,
\emph{The Malliavin Calculus and Related Topics},
2nd ed., Probability and Its Applications,
Springer-Verlag, Berlin, 2006.

\bibitem{NualartPardoux}
D.~Nualart and E.~Pardoux,
\emph{Stochastic calculus with anticipating integrands},
Probab. Theory Related Fields \textbf{78} (1988), no.~4, 535--581.

\bibitem{ReedSimonI}
M.~Reed and B.~Simon,
\emph{Methods of Modern Mathematical Physics. I: Functional Analysis},
revised and enlarged ed., Academic Press, New York, 1980.

\end{thebibliography}

\end{document}
